\section{Equilibrium in the Floor Economy}\label{sec:equilibrium}

\Cref{sec:incentives} compares the plans of one agent. This section aggregates those plans for the case in which the market is at the floor, and derives the loan book, the growth of the floor and the issuer's fee income. It is a conditional benchmark: arrivals and horizons are given, and \cref{sec:participation} shows that agents who value cash only would not supply them. The case is the one in which the price-taking assumption (A4) holds exactly, so that nothing is assumed about how agents interact.

\subsection{The economy}

The schedule is flat: every token is minted and redeemed at the floor price $\Pf(t)$, and capacity never binds. There is a continuum of holders. New holders arrive at a constant flow and buy $\nu>0$ tokens per unit of time, paying $(1+f_b)\Pf(t)$ per token. Each holder has a \emph{horizon} $T>0$, known to the holder on arrival. The distribution $F$ of horizons is over tokens bought, not over accounts: $\nu\,dF(T)$ is the flow of tokens bought by holders with horizon $T$. It has a density and a finite mean $\bar\tau$. If accounts differ in size, an account that buys $\ell$ tokens counts $\ell$ times, so that in terms of accounts the locked share \eqref{eq:U} below is $\mathbb E[\ell\,T\,\mathbf 1\{T>T_0\}]/\mathbb E[\ell\,T]$. At age $T$ the holder redeems all tokens, closing any loan first, and receives $(1-f_s)\Pf$ per token. Holders satisfy (A1) with $\rho>0$, (A2) and (A3), the policy LTV is a constant $\gamma\le\gbar$, and $\gbar+f_s\le1$.

The issuer applies a fraction $\alpha\in[0,1]$ of every fee to the floor in the transaction that pays it, and skims the rest. With a flat schedule, raising the floor by $d\Pf$ costs $S\,d\Pf$ (\cref{sec:raise}).

Because the schedule is flat, what a holder receives per token at exit is $(1-f_s)\Pf$ at the exit date, whatever the other holders do and whatever the holder's own borrowing was. A holder is infinitesimal and does not affect the path of the floor. Assumption (A4) therefore holds.

\begin{definition}[Stationary equilibrium with given arrivals]
A \emph{stationary equilibrium} consists of a growth rate $g\ge0$ of the floor, $\Pf(t)=\Pf(0)e^{gt}$, and a borrowing plan for each horizon $T$, such that (i) each plan maximizes the present value of its holder's cash flows given the path of the floor, and (ii) the floor grows at the rate $g$ when the supply is stationary and the fees are those generated by the plans.
\end{definition}

\subsection{Equilibrium}

Let $T_0=\frac1\rho\ln\frac1{1-\phi}$, the break-even horizon of \cref{prop:util}.

\begin{theorem}[Equilibrium of the floor economy]\label{thm:eq}
\begin{enumerate}[label=(\alph*)]
\item \emph{Plans.} For every continuous non-decreasing path of the floor, the following plan is optimal for a holder with horizon $T$: borrow every unit of borrowing capacity as soon as it becomes available while the remaining horizon exceeds $T_0$; borrow nothing during the last $T_0$ before exit; repay at exit. For $T\ne T_0$ it is the only optimal plan among those that repay at exit only, up to a set of ages of measure zero; a holder with $T=T_0$ on a constant floor is indifferent about borrowing at arrival. A plan that repays before exit is weakly dominated. If $\phi=0$ a holder can also repay and borrow again at the same moment at no cost, which adds to gross borrowing and changes no payment and no debt. The plan does not depend on the path of the floor or on the plans of other holders.
\item \emph{Loan book.} In the stationary state the supply is $S=\nu\bar\tau$, the locked supply is $L=\nu\,\mathbb E\bigl[T\,\mathbf 1\{T>T_0\}\bigr]$, and the share of the supply that is locked is
\begin{equation}\label{eq:U}
U=\frac{\mathbb E\bigl[T\,\mathbf 1\{T>T_0\}\bigr]}{\bar\tau}.
\end{equation}
\item \emph{Growth.} All stationary equilibria have the same growth rate and the same aggregate loan book; individual plans are determined for almost every horizon, apart from the costless round trips at $\phi=0$, which generate no fee. The growth rate is
\begin{equation}\label{eq:g}
g^\ast=\frac{\alpha\,\bigl[f_b+f_s+\phi\gamma\,(1-F(T_0))\bigr]}{\bar\tau-\alpha\,\phi\gamma\;\mathbb E\bigl[(T-T_0)^+\bigr]} .
\end{equation}
\end{enumerate}
\end{theorem}
\begin{proof}
A \emph{plan} of a holder is a non-decreasing right-continuous function of age giving cumulative borrowing, together with finitely many repayments before exit and the repayment of the remaining debt at exit; the debt may never exceed the limit. All plans have finite discounted cash flows: with $N$ repayments before exit, gross borrowing is at most $N+1$ times the limit at exit.

(a) We first show that a plan with a repayment before exit is weakly dominated by one without. Let a plan repay $r$ at age $a_1<T$, when its debt is $D_1$, and make no other repayment between $a_1$ and exit. Change it as follows: do not repay at $a_1$, and cancel later borrowings, in order, up to a total of $r$. The debt of the changed plan at any later age $a$ is the larger of $D_1$ and the debt of the original plan; both are at most the limit at age $a$, because the limit is non-decreasing, so the changed plan is feasible. It saves the payment $r$ at $a_1$; it forgoes at most $(1-\phi)c$ of later proceeds, where $c\le r$ is the amount canceled; and it repays $r-c$ more at exit. All of these later amounts are discounted more heavily than $a_1$, and $(1-\phi)c+(r-c)\le r$, so the present value does not fall. By (A2) the tokens are idle until exit in both plans. Applying this to each repayment in turn, starting with the last, we may restrict attention to plans that repay at exit only. Such a plan is a non-decreasing function $B(a)$ with $B(0^-)=0$, the cumulative amount borrowed up to age $a\in[0,T]$ (a loan at arrival is a jump at $a=0$), with $B(a)\le \bar B(a)$, where $\bar B(a)=\gamma\Pf(t_0+a)\ell$ is the borrowing limit at age $a$ of a holder who arrived at $t_0$ with $\ell$ tokens. A unit borrowed at age $a$ pays $1-\phi$ at once and costs $1$ at age $T$, so the present value of the plan, apart from terms that do not depend on $B$, is the Stieltjes integral $\int_{[0,T]} \omega(a)\,dB(a)$ with $\omega(a)=(1-\phi)e^{-\rho a}-e^{-\rho T}$. The function $\omega$ is strictly decreasing, and it is positive exactly for $a<a^\ast=T-T_0$. If $a^\ast\le0$ the integral is maximized by $B\equiv0$. Otherwise, integrating by parts on $[0,a^\ast]$, where $\omega(a^\ast)=0$,
\[
\int_0^{T}\omega\,dB=\int_0^{a^\ast}B(a)\,\bigl(-\omega'(a)\bigr)\,da+\int_{a^\ast}^{T}\omega\,dB .
\]
The first term is at most $\int_0^{a^\ast}\bar B\,(-\omega')\,da$, with equality if and only if $B=\bar B$ almost everywhere on $[0,a^\ast]$; the second is at most $0$, with equality if and only if $B$ is constant on $[a^\ast,T]$. This is the stated plan, and for $T\ne T_0$ it is the only maximizer up to ages of measure zero. If $T=T_0$ then $a^\ast=0$ and $\omega(0)=0$, so a loan at arrival neither adds nor subtracts value. At exit the debt is at most $\gamma\Pf\ell\le(1-f_s)\Pf\ell$, so the loan can be closed from the proceeds (\cref{prop:netclose}). Nothing in the argument uses the form of the path or the behavior of others.

(b) A token bought by a holder with horizon $T$ stays in the supply for a time $T$, so the stationary supply is $\nu\int T\,dF=\nu\bar\tau$. By (a), holders with $T>T_0$ lock their tokens on arrival and keep them locked until exit, and the others never lock; hence $L=\nu\int_{T_0}^\infty T\,dF$.

(c) Measure fee flows in units of the current floor price. Buyers pay $\nu f_b$ and redeemers $\nu f_s$ per unit of time. Origination fees are $\phi$ times the gross flow of new debt. By (a) that flow has two parts: the first loans of arriving holders with $T>T_0$, which are $\nu\gamma\,(1-F(T_0))$; and the top-ups of holders whose remaining horizon exceeds $T_0$, whose tokens number $\nu\,\mathbb E[(T-T_0)^+]$ and whose limit grows at the rate $g$, which are $g\gamma\nu\,\mathbb E[(T-T_0)^+]$. The floor rises at the rate $g$ when the fraction $\alpha$ of the fees equals the lift cost $gS=g\nu\bar\tau$:
\[
g\,\nu\bar\tau=\alpha\,\nu\Bigl[f_b+f_s+\phi\gamma\,(1-F(T_0))+\phi\gamma\,g\,\mathbb E\bigl[(T-T_0)^+\bigr]\Bigr].
\]
This equation is linear in $g$, and its coefficient $\bar\tau-\alpha\phi\gamma\,\mathbb E[(T-T_0)^+]$ is positive because $\alpha\phi\gamma<1$ and $\mathbb E[(T-T_0)^+]\le\bar\tau$. Its unique solution is \eqref{eq:g}. The plans in (a) do not depend on $g$, and because $F$ has a density the horizon $T_0$ has probability zero, so the indifference there does not affect the aggregates.
\end{proof}

Part~(a) is the reason the analysis of \cref{sec:incentives} was not misleading for being about one agent: at the floor, borrowing is a decision problem with a dominant plan, not a game. Credit is non-rivalrous (\cref{cor:nonrival}), and the only channel between holders is the floor, which every holder takes as given and which affects the two plans being compared in the same way.

For exponentially distributed horizons with mean $\tau$, \eqref{eq:U} is
\begin{equation}\label{eq:Uexp}
U=\Bigl(1+\frac{T_0}\tau\Bigr)e^{-T_0/\tau},
\end{equation}
and $\mathbb E[(T-T_0)^+]=\tau e^{-T_0/\tau}$. With $\rho=5\%$ a year and $\tau$ of one year, the locked share is $98.2\%$ for an origination fee of $1\%$ and $87.5\%$ for a fee of $3\%$. With $f_b=f_s=1\%$, $\phi=3\%$, $\gamma=0.9$ and all fees applied to the floor, the floor grows at $3.5\%$ a year (script block N11).

\begin{corollary}[Comparative statics]\label{cor:statics}
The locked share $U$ is non-increasing in the origination fee $\phi$ and non-decreasing in the outside rate $\rho$; it does not depend on $\gamma$, $f_b$, $f_s$ or $\alpha$. The growth rate $g^\ast$ is non-decreasing in $\alpha$, $f_b$, $f_s$ and $\gamma$.
\end{corollary}
\begin{proof}
$T_0$ is increasing in $\phi$ and decreasing in $\rho$, and $U$ is non-increasing in $T_0$. In \eqref{eq:g} the numerator is non-decreasing and the denominator non-increasing in each of $\alpha$ and $\gamma$, and $f_b$, $f_s$ enter the numerator only.
\end{proof}

\subsection{The issuer's fee}\label{sec:fee}

Suppose the issuer chooses the origination fee before holders arrive and skims all fees ($\alpha=0$), so that the floor is constant. The flow of origination fees is $\phi\gamma\nu\,(1-F(T_0(\phi)))$ in units of the floor price. This is a gross figure: it does not charge the issuer for the time during which the lent cash is out of the reserve. We treat the gross fee first and the issuer's present value second.

\begin{proposition}[Gross-fee optimum]\label{prop:fee}
With exponentially distributed horizons of mean $\tau$ and $\alpha=0$, the flow of origination fees is proportional to $\phi\,(1-\phi)^{1/(\rho\tau)}$ and is maximized at
\[
\phi^\ast=\frac{\rho\tau}{1+\rho\tau}.
\]
At $\phi^\ast$ the break-even horizon is $T_0=\tau\,\ln(1+\rho\tau)/(\rho\tau)$, and the locked share \eqref{eq:Uexp} tends to $2/e\approx74\%$ as $\rho\tau\to0$.
\end{proposition}
\begin{proof}
$1-F(T_0)=e^{-T_0/\tau}=(1-\phi)^{1/(\rho\tau)}$. The derivative of $\phi(1-\phi)^k$ with $k=1/(\rho\tau)$ is $(1-\phi)^{k-1}\,[1-\phi-k\phi]$, which changes sign from positive to negative at $\phi=1/(1+k)$. At that fee $1-\phi^\ast=1/(1+\rho\tau)$, so $T_0/\tau=\ln(1+\rho\tau)/(\rho\tau)\to1$.
\end{proof}

With $\rho=5\%$ and $\tau$ of one year, $\phi^\ast$ is $4.8\%$ and the locked share at that fee is $74.5\%$ (script block N11). To first order $\phi^\ast$ is the product of the outside rate and the mean horizon: it prices the mean holding period at the borrowers' own rate.

Now let the issuer value cash at a rate $r_I\ge0$, the return the reserve could earn if it were not lent. A loan of $b$ that is repaid after a time $T$ costs the issuer $(1-\phi)b$ today and returns $b$ at $T$. The loan carries no credit risk, because the locked tokens are redeemable for at least the debt, so $r_I$ is a pure time value.

\begin{proposition}[The issuer's side of a loan]\label{prop:issuer}
\begin{enumerate}[label=(\alph*)]
\item The issuer's present value of a loan of $b$ with horizon $T$ is $b\,[e^{-r_IT}-(1-\phi)]$. For $r_I>0$ it is positive if and only if $T<\frac1{r_I}\ln\frac1{1-\phi}$; for $r_I=0$ it equals the fee $\phi b$.
\item If $r_I=\rho$, the issuer's present value is the negative of the borrower's gain \eqref{eq:util}. Every loan from which a borrower gains then has negative value to the issuer, at every fee $\phi<1$; a loan taken at indifference has zero value to both.
\item With exponentially distributed horizons of mean $\tau$, $\alpha=0$ and $r_I<\rho$, the issuer's present value of lending per arriving token is
\[
\gamma\,(1-\phi)^{1/(\rho\tau)}\Bigl[\frac{(1-\phi)^{r_I/\rho}}{1+r_I\tau}-(1-\phi)\Bigr],
\]
and it is maximized at
\[
\phi_I=1-(1+\rho\tau)^{-\rho/(\rho-r_I)} .
\]
$\phi_I$ equals the gross-fee optimum $\phi^\ast$ if $r_I=0$, increases in $r_I$, and tends to $1$ as $r_I\to\rho$.
\end{enumerate}
\end{proposition}
\begin{proof}
(a) is the present value of the two payments. (b) The borrower's gain is $b\,[(1-\phi)-e^{-\rho T}]$, and the issuer's value is its negative. (c) A loan is taken at arrival by holders with $T>T_0$ and repaid at $T$; with a constant floor there are no top-ups. Per arriving token the value is $\gamma\int_{T_0}^\infty[e^{-r_IT}-(1-\phi)]\,\tau^{-1}e^{-T/\tau}\,dT=\gamma e^{-T_0/\tau}\,[e^{-r_IT_0}/(1+r_I\tau)-(1-\phi)]$, and $e^{-T_0/\tau}=(1-\phi)^{1/(\rho\tau)}$, $e^{-r_IT_0}=(1-\phi)^{r_I/\rho}$. Write $z=1-\phi$, $m=1/(\rho\tau)$ and $\theta=r_I/\rho<1$. The value is $\gamma\,[z^{m+\theta}/(1+r_I\tau)-z^{m+1}]$, whose derivative in $z$ vanishes where $z^{\theta-1}=(m+1)(1+r_I\tau)/(m+\theta)=1+\rho\tau$, because $m+\theta=(1+r_I\tau)/(\rho\tau)$ and $m+1=(1+\rho\tau)/(\rho\tau)$. The value is zero at $z=0$, negative at $z=1$ if $r_I>0$, and positive at the stationary point, which is therefore the maximum. Solving for $z$ gives $\phi_I$.
\end{proof}

Part (b) is the observation that a zero-interest loan is a transfer of time value, and that at a common discount rate the borrower gains exactly what the lender loses. Lending is worth something to both sides only if the issuer's cash is worth less to the issuer than to the borrower. The gross-fee optimum $\phi^\ast$ is the issuer's best fee in the extreme case $r_I=0$, in which the reserve has no other use; a reserve that sits idle by design is such a case. If the reserve could earn $3\%$ while borrowers discount at $5\%$ and $\tau$ is one year, the best fee is $\phi_I=11.5\%$, not $4.8\%$, and the locked share at that fee is $30\%$ (script block N13). A fee far below $\phi^\ast$, such as $1\%$, leaves the time value to borrowers and produces a loan book close to the whole supply.

\subsection{Participation}\label{sec:participation}

\Cref{thm:eq} takes the arrivals as given. We now ask whether a holder with preferences (A1) and (A2) would choose to arrive. Take a stationary equilibrium with growth rate $g$, and measure a holder's cash flows per token in units of the floor price on arrival. A holder with horizon $T$ pays $1+f_b$ at age $0$ and receives $(1-f_s)e^{gT}$ at age $T$. If $T>T_0$ the holder also receives $(1-\phi)\gamma$ at age $0$, receives $(1-\phi)\gamma g e^{ga}\,da$ at each age $a\in(0,T-T_0)$ from top-ups, and repays $\gamma e^{g(T-T_0)}$ at age $T$. The \emph{entry value} $V(T)$ is the present value of these flows at the rate $\rho$:
\begin{multline}\label{eq:V}
V(T)=-(1+f_b)+(1-f_s)e^{(g-\rho)T}
+\mathbf 1\{T>T_0\}\,\gamma\Bigl[(1-\phi)-e^{-\rho T}\\
+g\int_0^{T-T_0}\bigl((1-\phi)e^{(g-\rho)a}-e^{-\rho T}e^{ga}\bigr)\,da\Bigr].
\end{multline}
Let $\bar V=\int V\,dF$ be its mean over arriving tokens. Three stationary quantities enter the result, all in units of the current floor price. The debt is $D_n$. The issuer's cash is $A_n=\nu\bar\tau-D_n$: we take a state without retained surplus, $A+D=R(S)$, and all fees are applied to the floor or skimmed as they are paid, so the reserve continues to hold exactly $R(S)-D$. The skim is $s$, the fraction $1-\alpha$ of the fee flow. Finally let $X$ be the \emph{continuation value} of the holders present at a given date: the present value at the rate $\rho$ of all payments they will receive less all payments they will make from that date on.

\begin{theorem}[Entry value]\label{thm:entry}
In a stationary equilibrium without retained surplus, with $g<\rho$ and $\gamma+f_s\le1$,
\begin{equation}\label{eq:entry}
\nu\,\bar V=-\bigl(g\,A_n+s\bigr)-(\rho-g)\,X,
\end{equation}
with $A_n\ge(1-\gamma)\nu\bar\tau\ge0$, $s\ge0$ and $X\ge0$. Hence $\bar V\le0$, and $\bar V<0$ if $\gamma+f_s<1$.
\end{theorem}
\begin{proof}
Let $N$ be the net flow of payments from the issuer to all holders, in units of the current floor price. It is constant in a stationary state. The issuer's cash $A_n\Pf(t)$ changes by what holders pay in net, less the skim: $\frac{d}{dt}\bigl(A_n\Pf\bigr)=-N\Pf-s\Pf$, so $N=-(gA_n+s)$. The debt of a holder is at most $\gamma$ times the current floor value of the holder's tokens, so $D_n\le\gamma\nu\bar\tau$ and $A_n\ge(1-\gamma)\nu\bar\tau$.

Fix a date $0$ and take present values at that date, in units of $\Pf(0)$. All payments between the issuer and holders from date $0$ on are worth $\int_0^\infty e^{-\rho t}Ne^{gt}\,dt=N/(\rho-g)$. Each of these payments goes to or comes from a holder who either is present at date $0$ or arrives later. Those of the first kind are worth $X$. A holder who arrives at $t>0$ with horizon $T$ has cash flows worth $V(T)\Pf(t)$ per token at date $t$, so the payments of the second kind are worth $\int_0^\infty e^{-\rho t}\nu\bar Ve^{gt}\,dt=\nu\bar V/(\rho-g)$. The integrals converge absolutely because $g<\rho$ and every term of \eqref{eq:V} is bounded in $T$. Hence $N/(\rho-g)=X+\nu\bar V/(\rho-g)$, which is \eqref{eq:entry}.

For the sign of $X$, take a holder at some age, with $r$ of the horizon remaining, tokens worth $P$ at the current floor and $P'\ge P$ at the floor on exit, and current debt $d\le\gamma P$. The holder's remaining flows are the redemption $(1-f_s)P'$ and the repayment of $d$ at exit, worth $e^{-\rho r}[(1-f_s)P'-d]\ge e^{-\rho r}(1-f_s-\gamma)P'\ge0$, and the future top-ups with their repayments, each of which has a non-negative present value because it is borrowed while $\omega>0$ in the proof of \cref{thm:eq}. If $\gamma+f_s<1$ the first amount is positive for every holder, so $X>0$.
\end{proof}

The identity is a budget identity, and we do not present it as more: what entrants as a group receive is what the mechanism pays out, and a mechanism that accumulates or skims cash pays out less than it takes in. Its content is the size and the sign of the terms. It has a direct reading. Entrants as a group pay for two things. The first term is the cash that the issuer accumulates or skims: a floor that rises at the rate $g$ on a stationary supply needs a reserve that grows at the rate $g$, and only holders can supply it. The second is the time cost of the capital that holders keep in the system, which appreciates at $g$ while its owners discount at $\rho$. The loans return part of that cost to borrowers, and \eqref{eq:entry} shows that they cannot return all of it.

\begin{corollary}[No voluntary entry of cash-only agents]\label{cor:noentry}
Let potential buyers satisfy (A1) and (A2) and buy only if their entry value is non-negative, whatever the distribution of horizons among them. Then there is no stationary equilibrium of the floor economy in which a positive flow of them buys, the horizons of those who buy have a density and a finite mean, $g<\rho$ and $\gamma+f_s<1$.
\end{corollary}
\begin{proof}
In such an equilibrium the horizons of the agents who buy have a distribution $F$ and the floor grows at the rate that \cref{thm:eq} gives for $F$. \Cref{thm:entry} applies to that pair, however the set of buyers came about: the mean of $V$ over buyers is negative, so a set of buyers of positive measure has $V(T)<0$, which contradicts their decision to buy.
\end{proof}

Under these conditions the mechanism cannot carry itself on the motive that \cref{sec:incentives} gives its borrowers. The conditions matter: with $\gamma=1$ and no fees the entry value is zero for every horizon, and the case $g\ge\rho$ is discussed below. Buyers must receive something that the cash flows above do not contain, worth at least $-\bar V$ per token on average: a price above the floor that they expect to rise, a payment on held or locked tokens, or exposure that they value for its own sake. In the example after \cref{thm:eq} ($\rho=5\%$, $\tau$ of one year, $\phi=3\%$, $\gamma=0.9$, $f_b=f_s=1\%$, $\alpha=1$), $g^\ast=3.52\%$ and the mean entry value is $-1.10\%$ of the floor price on arrival. A holder with a horizon of one year has $V=-1.75\%$: the purchase and redemption cost $3.45\%$, and the loan and its top-ups return $1.70\%$. The entry value is negative for horizons below $1.61$ years, which is $80\%$ of arriving tokens. In \eqref{eq:entry}, per unit of $\nu$, $A_n=0.216$ and $X=0.234$, and the two terms are $0.76\%$ and $0.35\%$. A benefit paid at a constant rate of $1.16\%$ of the arrival floor price per token and year for as long as a token is held would bring the mean entry value to zero (script block N12).

The theorem does not cover $g\ge\rho$. By \eqref{eq:g} that case arises when horizons are short relative to the fees, roughly when $\alpha\,(f_b+f_s+\phi\gamma)$ exceeds $\rho\bar\tau$. The absolute-convergence argument used in the proof then fails, so this proof does not establish the identity in that regime; it can still hold there (at $g=\rho$ it reduces to $\nu\bar V=-(gA_n+s)$ whenever the quantities are finite), and the integrability and boundary conditions need a separate analysis. Any gain of a long-horizon holder must be funded by other cash flows, including later buyers. Purchases measured in the reserve asset, $\nu\Pf(t)$, must grow at a rate of at least $\rho$ without end. The situation resembles the one studied in the literature on rational bubbles, where an asset can be valued above its payouts only if the economy that buys it grows fast enough \citep{tirole1985,santos1997}. We note the resemblance only: we have not derived a fundamental value, a bubble component or a participation equilibrium for this regime, and we do not analyze it.

\subsection{Service value and refundable capital}\label{sec:deposit-benchmark}

The participation result concerns the cash flows in \eqref{eq:V}. A service can make a token useful without paying a cash dividend, but a comparison must hold that service fixed. The following elementary benchmark states the capital cost that such a benefit would have to cover; it does not supply an adoption equilibrium.

\begin{samepage}
\begin{proposition}[Refundable-capital benchmark]\label{prop:deposit-benchmark}
Let the mint and redemption price be a constant $P>0$, with no fees or interest. A holder buys one token, immediately borrows $\gamma P$ against it, and closes and redeems at a fixed time $T>0$, where $0\le\gamma\le1$ and $\gamma=0$ denotes no borrowing. Suppose the holder discounts at $\rho>0$ and receives service benefit at a constant rate $\sigma\ge0$, in reserve-equivalent units, whether the token is free or locked. The value of entering is
\begin{equation}\label{eq:deposit-benchmark}
V_{\mathrm{service}}(T)
=\left(\frac{\sigma}{\rho}-(1-\gamma)P\right)(1-e^{-\rho T}).
\end{equation}
It is non-negative if and only if $\sigma\ge\rho(1-\gamma)P$. The dated net cash flows and service benefits are identical to those of a refundable deposit of $(1-\gamma)P$ that provides the same service until $T$.
\end{proposition}
\end{samepage}
\begin{proof}
The initial net payment is $-(1-\gamma)P$ and the terminal net receipt is $(1-\gamma)P$. The service is worth $\int_0^T \sigma e^{-\rho t}\,dt=\sigma(1-e^{-\rho T})/\rho$. Adding the three terms gives \eqref{eq:deposit-benchmark}. A deposit of the stated amount with the same refund date and service has precisely these flows.
\end{proof}

For example, with $P=1{,}000$ and $\rho=5\%$ a year, the required annual service benefit is $50$ without borrowing and $5$ at $\gamma=0.9$. Credit reduces the net capital commitment, but a deposit of $100$ replicates the latter arrangement. The comparison concerns net dated flows: without atomic purchase-and-borrow settlement, the token route also needs temporary funding of the gross purchase. Fees, a changing floor and purchases above the floor change the comparison. A service benefit is not cash available to pay the platform's operating costs.

The economic question is therefore which friction, commitment or service makes this arrangement preferable to a refundable deposit or outside finance. Token adoption, productive platform finance and coordination through costly purchases already have developed models \citep{cong2021tokenomics,cong2022finance,limann2025}. Reserve credit changes both the capital committed by a user and the user's exit rights. Establishing whether those changes support participation requires a model of those frictions and the platform's budget, not the addition of an unexplained positive utility term to \eqref{eq:V}.

\subsection{What each assumption supplies}

\Cref{tab:assumptions} states, for each assumption of \cref{sec:incentives}, the results that rest on it and what replaces them when it fails.

\begin{table}[t]
\centering\small
\caption{The four assumptions: what each supplies and what changes when it fails.}\label{tab:assumptions}
\begin{tabular}{@{}p{0.17\linewidth}p{0.36\linewidth}p{0.41\linewidth}@{}}
\toprule
Assumption & Supplies & If it fails \\
\midrule
(A1) cash only, outside rate $\rho$ & The break-even horizon $T_0$, the value of deferral, the plans of \cref{thm:eq}. & The model has no price above the floor and no beliefs about it. An agent who expects that price to rise may borrow to buy more tokens while valuing cash only; what is missing is portfolio choice under beliefs, not a different preference. Under (A1), (A2), $g<\rho$ and $\gamma+f_s<1$, buyers do not arrive (\cref{cor:noentry}). \\
(A2) tokens pay nothing & Locking is free, so borrowing and holding differ only in the loan. & If locked tokens earn at least what free tokens earn, the results hold with a larger gain (\cref{prop:general}). If only free tokens earn, locking has a cost and $T_0$ rises. \\
(A3) no interest, no maturity, proportional sell fee & Deferral of repayment is free; the debt at exit is known at origination. & With interest, deferral has a price and repayment before exit can be optimal; debt can outgrow the floor value unless accrual is capped (\cref{sec:model}). \\
(A4) price taker & One agent's plan does not change the agent's own proceeds. Exact on a flat schedule. & Above the floor, proceeds depend on the order of exit (\cref{cor:run}) and on other agents' trades (\cref{sec:extraction}); exit becomes a game, which is not solved here. \\
\bottomrule
\end{tabular}
\end{table}

\subsection{What the floor economy leaves out}

Three limits apply. First, horizons and arrivals are exogenous. \Cref{cor:noentry} shows that this is a substantive restriction and not a simplification: the model has no demand side, and it says nothing about how fees affect demand. Second, above the floor the curve return of a redemption depends on the order of exit (\cref{cor:run}), which makes holders' exit decisions interdependent; the equilibrium of that game is not solved here, and \cref{thm:eq} should not be read as covering it. Third, borrowers are assumed to want cash. A borrower who uses the loan to buy more tokens is taking a leveraged position on the price. The model does not represent that use, and in the Baseline data it accounts for \ELeverAmtPct\% of new debt by amount (\cref{sec:empirics}). The data do not show what these buyers expect; the deployed design also pays rewards, and beliefs are not observed. Price exposure, fee rewards and service value are distinct candidate benefits; the data do not separate their contributions to entry. For this reserve-lending contract, we leave a model of participation and leverage under stated beliefs, with rewards and a funded service where relevant, to future work.
